% !TeX root = ../papers/nc-p3.tex
\section{重根の二位置交換の固定点とEnriques商}
\label{sec:kummer-exchange-quotient}
\begin{theorem}\label{thm:kummer-exchange-enriques}
通常倍率の各幾何点で、規格化した二行列の交換は判別Kummer曲面の最小解消上の
固定点のない対合となる。従ってこの交換による解消曲面の商はEnriques曲面である。
これは印と倍率を固定した表示の二位置交換に関する商である。
\end{theorem}
\begin{proof}
通常倍率の連結な基底を$B$とする。曲面族の相対勾配の基点は互いに交わらない
十六の節点切断である。節点切断に沿って相対Hesse行列が非退化なので、アフィン
座標の一次偏微分は切断の余接空間の基底を与え、Nakayamaの補題で切断のイデアルを生成する。
従って極写像のグラフ閉包$Y$はこれら十六切断の同時吹上げである。
各例外曲線は滑らかな接錐二次曲線である。吹上げの各アフィン座標では、
二次方程式の非特異な射影接錐により$Y\to B$は例外上でも滑らかである。
これは固有な曲面族で、各ファイバーは通常二重点の最小解消である。
極写像が対合なので、二つの射影因子の交換は$Y$の正則な自己同型$\tau$となる。

$N=[u]$においてグラフに固定点があれば
\[
 F(u)=0,\qquad u^t\operatorname{adj}(T)u=0,\qquad
 u\wedge T\nabla F(u)=0 .
\]
滑らかな点ではEuler公式$u^t\nabla F=4F$を使う。
節点では極写像の極限は$u$を通る平面なので、グラフ上の第二座標は
$u^tT^{-1}v=0$を満たす。従って必要条件は節点にも及ぶ。
$(z,\lambda)=(3,2)$の整数係数でこの斉次閉集合を素数1009へ還元する。
先頭の非零座標を1とする四つの部分で、最初の三つではGröbner基底が$(1)$、
最後では二次式が非零定数である。
\path{computations/verify_kummer_exchange_quotient.py}がこれらを検算する。
射影閉集合の固有性により標数零のこの有理ファイバーにも幾何的な解はない。

固定点スキーム$Y^\tau$は$B$上固有かつ滑らかである。後者を詳しく見ると、
平方零のイデアルを持つ試験環への固定点の持ち上げについて、$Y/B$の滑らかさで
まず一つの持ち上げを取れる。持ち上げ全体は接ベクトルを係数とするアフィン空間で、
対合の作用はアフィン線型である。持ち上げとその対合像の中点を取れば固定した
持ち上げを得る。これは形式的滑らかさを示し、有限表示なので滑らかさとなる。
従って像は開かつ閉であり、連結基底の一点で空だから全体で空である。

解消四次曲面はK3曲面である。実際、四次超曲面の標準束は自明で通常二重点は
クレパントな有理特異点であり、四次の完全列から$H^1(\mathcal O)=0$である。
自由な対合による商は滑らかで射影的、被覆は次数2のエタール射である。
従って商でも$H^1(\mathcal O)=0$、$\chi(\mathcal O)=1$で、標準束の二乗は
ノルムにより自明である。Serre双対性から標準束自身は切断を持たず、自明でない。
これはEnriques曲面の条件を満たす。他の文字置換も許した添字保持同型の判定には、
別途四位置の作用を考慮する必要がある。
\end{proof}
